\documentclass{article}
\usepackage[margin=1in]{geometry}
\usepackage{hyperref}

\title{Problem Set 4}
\author{Sumedha Ray}
\date{February 24, 2026}

\begin{document}
\maketitle

\section*{Data Sources I Would Like to Scrape}

There are several data sources I would be interested in scraping for research purposes.
First, the NBER Working Papers website (\url{https://www.nber.org/papers}) provides 
access to working papers in economics, which could be scraped to analyze 
trends in research topics, citations, and author networks over time. Second, 
EconJobMarket.org (\url{https://econjobmarket.org}) contains job market candidate 
profiles and placement outcomes, which could be used to study patterns in PhD placement 
across institutions, fields, and demographic groups. Third, the IZA Discussion Papers 
database (\url{https://docs.iza.org}) which has a large collection of labor economics 
papers that could be scraped to study research trends in gender economics and 
discrimination. 
I would also be interested in scraping the World Bank Open Data API provides extensive development indicators for countries over time. 


\section*{Q5(d): Object Types}
The object \texttt{mydf} is of class \texttt{tbl\_df}, \texttt{tbl}, and 
\texttt{data.frame} -- that is, it is a tibble/dataframe. The object 
\texttt{mydf\$date} is of class \texttt{character}, meaning it is stored as 
a character/string vector rather than a numeric or date type.

\section*{Q6(7): Spark vs. Tibble Classes}
\texttt{df1} (a tibble) has class \texttt{tbl\_df}, \texttt{tbl}, \texttt{data.frame}.
\texttt{df} (a Spark dataframe) has class \texttt{tbl\_spark}, \texttt{tbl\_sql}, 
\texttt{tbl\_lazy}, \texttt{tbl}.

\section*{Q6(8): Column Names}
The column names differ between the tibble and the Spark dataframe. In the tibble,
column names use dots (e.g., \texttt{Sepal.Length}), while in Spark they use 
underscores (e.g., \texttt{Sepal\_Length}). This is because Spark/SQL does not 
allow dots in column names, so \texttt{sparklyr} automatically converts them.

\section*{Q6: Note on sparklyr}
Despite attempting to install \texttt{sparklyr} on OSCER, the installation was 
unsuccessful due to dependency failures, particularly with the \texttt{dbplyr} 
package. This appears to be a compatibility issue between the version of R 
available on OSCER (R 4.0.2) and the current version of \texttt{sparklyr}. 
As noted in the problem set instructions, students would not be penalized if 
\texttt{sparklyr} does not work on the first try.

\end{document}